% Relativistic Chapter VII, §67: Stability of inviscid relativistic Couette flow
% Relativistic generalization of Chandrasekhar, Chapter VII §67 (pp.~319--753)
% Conventions: see RELATIVISTIC_CONVENTIONS.md

\section{Analytical discussion of the stability of inviscid relativistic
  Couette flow}
\label{sec:rel-7-67}

We now consider the stability of a stationary, purely azimuthal,
relativistic flow between two coaxial cylinders of radii $R_1$ and~$R_2$.
In cylindrical coordinates $(r,\theta,z)$ the background four-velocity
is
\begin{equation}
\label{eq:rel-7-33}
u^\mu = \Lf\bigl(c,\,0,\,V(r),\,0\bigr),\qquad
V(r) = r\,\Omega(r),\qquad
\Lf = \Bigl(1-\frac{V^2}{c^2}\Bigr)^{-1/2},
\end{equation}
where $\Omega(r)$ is an arbitrary angular-velocity profile and
$\Lf$ denotes the Lorentz factor of the background flow.
The flow is sustained by a radial pressure gradient that balances
the relativistic centrifugal force (cf.\ the relativistic Euler
equation projected onto the radial direction):
\begin{equation}
\label{eq:rel-7-33b}
\frac{1}{\enthalpy}\frac{dp}{dr}
= \Lf^{2}\,\frac{V^{2}}{r},
\end{equation}
where $\enthalpy = (\varepsilon + p)/c^{2}$ is the relativistic
enthalpy density.  In the limit $V/c\to0$ (i.e.\ $\Lf\to1$,
$\enthalpy\to\rho$) this reduces to the classical balance
$\rho^{-1}dp/dr = V^{2}/r$.

Throughout this section the fluid is taken to be a \emph{perfect}
(inviscid, non-heat-conducting) relativistic fluid whose equation of
state is barotropic, $p = p(\varepsilon)$, with adiabatic sound speed
$\cs^{2}=(\partial p/\partial\varepsilon)_s\le c^{2}$.
Because the flow is incompressible in the sense that the
baryon density is uniform (the classical Couette assumption),
we work with an effective incompressible relativistic formulation
in which the enthalpy density $\enthalpy$ plays the role of the
inertial mass density.

%% ----------------------------------------------------------------
\subsection*{(a) \emph{Lagrangian displacement formulation in special
  relativity}}
\label{sec:rel-7-67a}

Consider an infinitesimal perturbation of~\eqref{eq:rel-7-33}.
Let the Eulerian perturbations of the three-velocity components be
$\delta v_r$, $\delta v_\theta$, $\delta v_z$, and write
$\varpi = \delta p/(\Lf^{2}\enthalpy)$ for the relativistic
pressure perturbation (the classical limit gives $\varpi\to\delta p/\rho$
as required).  Linearising the relativistic Euler equations about the
background flow, one obtains the perturbation equations in the
co-rotating frame.  Because the background depends only on~$r$,
we decompose every perturbation quantity in normal modes
\begin{equation}
\label{eq:rel-7-39}
\propto \exp\!\bigl[i(\sigma t + m\theta + kz)\bigr],
\end{equation}
where $m$ is an integer and $k$ is a real axial wave number.  The
Doppler-shifted frequency is
\begin{equation}
\label{eq:rel-7-44}
\sigma(r) = p + m\,\Omega(r),
\end{equation}
with $p$ the frequency in the laboratory frame.

In full analogy with the classical treatment, we introduce the
\emph{relativistic Lagrangian displacement} $\boldsymbol{\xi}$ through
\begin{equation}
\label{eq:rel-7-45}
\delta v_r = i\sigma\,\xi_r,\qquad
\delta v_\theta = i\sigma\,\xi_\theta
  - r\frac{d\Omega}{dr}\,\xi_r,\qquad
\delta v_z = i\sigma\,\xi_z.
\end{equation}
This definition ensures that $\boldsymbol{\xi}$ describes the actual
fluid-element displacement to first order in the perturbation
amplitude.  The relativistic solenoidal condition
$\nabla_\mu((\enthalpy/c^{2})\,\delta u^\mu)=0$, when projected onto the
spatial sector, yields
\begin{equation}
\label{eq:rel-7-46}
\frac{d\xi_r}{dr} + \frac{\xi_r}{r}
+ \frac{im}{r}\,\xi_\theta + ik\,\xi_z = 0.
\end{equation}

In terms of $\boldsymbol{\xi}$ the linearised relativistic momentum
equations become
\begin{equation}
\label{eq:rel-7-47}
\bigl[\sigma^{2}-\Phi_{\mathrm{rel}}(r)\bigr]\,\xi_r
= \frac{d\varpi}{dr}
  + \frac{2m\,\Omega}{r\,\sigma}\,\varpi,
\end{equation}
\begin{equation}
\label{eq:rel-7-48}
\sigma^{2}\,\xi_\theta - 2i\,\Omega\,\sigma\,\xi_r
= \frac{im}{r}\,\varpi,
\end{equation}
\begin{equation}
\label{eq:rel-7-49}
\sigma^{2}\,\xi_z = ik\,\varpi,
\end{equation}
where $\Phi_{\mathrm{rel}}(r)$ is the \emph{relativistic Rayleigh
discriminant}:
\begin{equation}
\label{eq:rel-7-53}
\boxed{%
\Phi_{\mathrm{rel}}(r)
= \frac{2\,\Lf^{4}\,\Omega}{r}\,
  \frac{d}{dr}\!\bigl(r^{2}\,\Lf^{2}\,\Omega\bigr).}
\end{equation}
In the non-relativistic limit $\Lf\to1$ this reduces to the
classical Rayleigh discriminant
$\Phi(r)=(2\Omega/r)\,d(r^{2}\Omega)/dr$.

\medskip
\noindent\textbf{Remark (relativistic corrections in $\Phi_{\mathrm{rel}}$).}
The additional $\Lf$-factors encode two distinct relativistic effects:
(i) the increase of inertia with velocity ($\Lf^{2}$ per factor of
$\enthalpy$), and (ii) the modification of angular-momentum transport
due to the non-linear velocity-addition law.  Writing
$\ell_{\mathrm{rel}}=\Lf^{2}\,r^{2}\,\Omega$ for the specific
relativistic angular momentum, one recognises
$\Phi_{\mathrm{rel}} = (2\Lf^{2}\Omega/r)\,d\ell_{\mathrm{rel}}/dr$.

Combining equations~\eqref{eq:rel-7-46}--\eqref{eq:rel-7-49} in
the same manner as in the classical analysis, one derives the
companion equation
\begin{equation}
\label{eq:rel-7-54}
\frac{1}{r}\frac{d}{dr}(r\,\xi_r)
- \frac{2m\,\Omega}{\sigma\,r}\,\xi_r
= \frac{1}{\sigma^{2}}
  \!\left(\frac{m^{2}}{r^{2}}+k^{2}\right)\varpi.
\end{equation}
Equations \eqref{eq:rel-7-47} and \eqref{eq:rel-7-54}, together with
the boundary conditions
\begin{equation}
\label{eq:rel-7-55}
\xi_r = 0 \quad\text{for } r = R_1 \text{ and } R_2,
\end{equation}
constitute a self-adjoint system for real~$\sigma$, exactly as in the
non-relativistic case.  Their solution can therefore be cast as a
variational problem---the entire merit of introducing the Lagrangian
displacement.

%% ================================================================
\subsection*{(b) \emph{The case $m=0$ (axisymmetric perturbations):
  the relativistic Rayleigh criterion}}
\label{sec:rel-7-67b}

When $m=0$ we have $\sigma = p$ and equations~\eqref{eq:rel-7-47},
\eqref{eq:rel-7-54} simplify to
\begin{equation}
\label{eq:rel-7-56}
\bigl[p^{2}-\Phi_{\mathrm{rel}}(r)\bigr]\,\xi_r
= \frac{d\varpi}{dr},
\end{equation}
\begin{equation}
\label{eq:rel-7-57}
\frac{1}{r}\frac{d}{dr}(r\,\xi_r)
= \frac{k^{2}}{p^{2}}\,\varpi.
\end{equation}
Eliminating $\varpi$ yields the relativistic Sturm--Liouville equation
\begin{equation}
\label{eq:rel-7-58}
\frac{d}{dr}\!\left\{
\frac{1}{r}\frac{d}{dr}(r\,\xi_r)\right\}
- k^{2}\,\xi_r
= -\frac{k^{2}}{p^{2}}\,\Phi_{\mathrm{rel}}(r)\,\xi_r.
\end{equation}
By the standard Sturm--Liouville theory applied to the boundary
conditions~\eqref{eq:rel-7-55} one deduces:

\emph{The characteristic values of $k^{2}\!/p^{2}$ are all positive if
$\Phi_{\mathrm{rel}}(r)>0$ throughout $(R_1,R_2)$, all negative if
$\Phi_{\mathrm{rel}}(r)<0$ throughout, and of both signs if
$\Phi_{\mathrm{rel}}$ changes sign in the interval.}

Multiplying \eqref{eq:rel-7-58} by $r\,\xi_r$ and integrating, one
obtains the relativistic variational expression
\begin{equation}
\label{eq:rel-7-60}
\frac{p^{2}}{k^{2}}
= \frac{\displaystyle\int_{R_1}^{R_2}
    \Phi_{\mathrm{rel}}(r)\,r\,\xi_r^{2}\,dr}
  {\displaystyle\int_{R_1}^{R_2}\!\left\{
    \frac{1}{r}\!\left[\frac{d}{dr}(r\,\xi_r)\right]^{2}
    + k^{2}\,r\,\xi_r^{2}\right\}dr}.
\end{equation}
Since the denominator is manifestly positive, the sign of $p^{2}$ is
governed entirely by the sign of
\begin{equation}
\label{eq:rel-7-61}
I_1 = \int_{R_1}^{R_2}
  \Phi_{\mathrm{rel}}(r)\,r\,\xi_r^{2}\,dr.
\end{equation}

\begin{quote}
\textbf{Relativistic Rayleigh criterion.}\quad
\emph{A necessary and sufficient condition for the stability of an
inviscid, relativistic, axisymmetric Couette flow against
axisymmetric perturbations is that the relativistic specific angular
momentum $\ell_{\mathrm{rel}}(r)=\Lf^{2}\,r^{2}\,\Omega$ be a
non-decreasing function of~$r$ throughout the flow domain:}
\begin{equation}
\label{eq:rel-rayleigh}
\boxed{%
\frac{d}{dr}\!\bigl(\Lf^{2}\,r^{2}\,\Omega\bigr)\ge 0
\quad\text{for all } r\in[R_1,R_2].}
\end{equation}
\emph{Equivalently, $\Phi_{\mathrm{rel}}(r)\ge 0$ everywhere.
If $\Phi_{\mathrm{rel}}<0$ anywhere, the flow is unstable.}
\end{quote}

\noindent\textbf{Non-relativistic limit.}\quad
When $V/c\to0$ we have $\Lf\to1$ and the criterion~\eqref{eq:rel-rayleigh}
reduces to the classical Rayleigh criterion
$d(r^{2}\Omega)/dr\ge0$.

\noindent\textbf{Physical interpretation.}\quad
In relativity the ``effective angular momentum'' is not $r^{2}\Omega$
but $\Lf^{2}r^{2}\Omega$, reflecting the relativistic increase of
inertia.  A flow that is marginally stable classically can therefore
become \emph{unstable} at relativistic velocities if the Lorentz-factor
gradient is sufficiently steep and negative.  Conversely, a classically
unstable arrangement can be \emph{stabilised} if $\Lf^{2}$ increases
outward rapidly enough.

%% ================================================================
\subsection*{(c) \emph{The case $m\neq0$ (non-axisymmetric
  perturbations): relativistic Howard semicircle theorem}}
\label{sec:rel-7-67c}

For $m\neq 0$ the problem is treated, following the classical analysis,
as a characteristic-value problem for~$k^{2}$ at a given (possibly
complex) frequency~$p$.  Equations~\eqref{eq:rel-7-47}
and~\eqref{eq:rel-7-54}, together with the boundary
conditions~\eqref{eq:rel-7-55}, form a Hermitian system when $p$ is
real; the characteristic values of~$k^{2}$ are therefore real for real~$p$.

The variational expression~\eqref{eq:rel-7-60} generalises to
\begin{equation}
\label{eq:rel-7-66}
k^{2}
= \frac{\displaystyle\int_{R_1}^{R_2}
    r\!\left[(\Phi_{\mathrm{rel}}-\sigma^{2})\,|\xi_r|^{2}
    - \frac{m^{2}|\varpi|^{2}}{r^{2}\sigma^{2}}\right]dr}
  {\displaystyle\int_{R_1}^{R_2}
    r\,|\varpi|^{2}/\sigma^{2}\,dr},
\end{equation}
with the subsidiary relation~\eqref{eq:rel-7-47} determining
$\xi_r$ in terms of~$\varpi$.

Using the standard procedure of Howard (1961), adapted to the
relativistic discriminant, one derives the following bound on unstable
modes.

\begin{quote}
\textbf{Relativistic Howard semicircle theorem.}\quad
\emph{Let $p = p_r + i\,p_i$ with $p_i>0$ (instability).
Then the complex frequency lies inside the semicircle}
\begin{equation}
\label{eq:rel-howard}
\boxed{%
\bigl(p_r + m\,\Omega_{\min}\bigr)
\bigl(p_r + m\,\Omega_{\max}\bigr)
+ p_i^{2}\;\le\;0,}
\end{equation}
\emph{where $\Omega_{\min}$ and $\Omega_{\max}$ are the minimum and
maximum values of $\Omega(r)$ in $[R_1,R_2]$.}
\end{quote}

\noindent\emph{Proof (sketch).}\quad
Multiply the governing equation by $r\,\xi_r^{*}/\sigma$ and integrate.
The imaginary part of the resulting integral identity is
\begin{equation}
\label{eq:rel-howard-imag}
p_i\!\int_{R_1}^{R_2}
  \frac{r\,|\xi_r|^{2}}{|\sigma|^{2}}
  \bigl[\Phi_{\mathrm{rel}}(r)-|\sigma|^{2}\bigr]\,dr = 0.
\end{equation}
For $p_i\neq0$ the integral must vanish, which requires $|\sigma|^{2}$
to lie between the extreme values of $\Phi_{\mathrm{rel}}$.
Since $\Phi_{\mathrm{rel}}\le
4\Lf^{4}\Omega^{2}\!\bigl[1+O(r\Omega'/\Omega)\bigr]$, one can bound
$|p+m\Omega|^{2}$ from above and below, yielding the
semicircle~\eqref{eq:rel-howard}.  The relativistic corrections enter
only through the modified discriminant and leave the algebraic structure
of Howard's argument intact.  In the limit $\Lf\to1$ one recovers the
classical Howard semicircle.\qed

\noindent\textbf{Consequence.}\quad
The range of unstable real frequencies $p_r$ is confined to the interval
$[-m\,\Omega_{\max},\,-m\,\Omega_{\min}]$, and the growth rate $p_i$
is bounded by
\begin{equation}
\label{eq:rel-howard-bound}
p_i \le \frac{|m|}{2}\,(\Omega_{\max}-\Omega_{\min}).
\end{equation}
This is the same bound as classically; the relativistic
effects do not enlarge the unstable semicircle (they modify
\emph{which} profiles are unstable via $\Phi_{\mathrm{rel}}$, but the
frequency bounds remain controlled by $\Omega_{\max}-\Omega_{\min}$).

%% ================================================================
\subsection*{(d) \emph{Energy principle for relativistic rotating flow}}
\label{sec:rel-7-67d}

The self-adjoint character of the system
\eqref{eq:rel-7-47}--\eqref{eq:rel-7-54} permits a formulation as an
energy principle.  Define the quadratic form
\begin{equation}
\label{eq:rel-7-energy}
\delta^{2}W[\boldsymbol{\xi}]
= \int_{R_1}^{R_2}\!\bigl\{
  \Lf^{2}\,\enthalpy\,\Phi_{\mathrm{rel}}\,r\,|\xi_r|^{2}
  + \Lf^{2}\,\enthalpy\,r\,\sigma^{2}
    \bigl(|\xi_\theta|^{2}+|\xi_z|^{2}\bigr)
  + r\,|\delta p|^{2}/(\Lf^{2}\,\enthalpy\,\cs^{2})
\bigr\}\,dr,
\end{equation}
where the three terms represent, respectively, the potential energy
associated with the angular-momentum stratification, the kinetic
energy of the perturbation, and the acoustic compression energy.
(In the incompressible limit $\cs\to\infty$ the last term vanishes.)

The system is linearly stable if and only if $\delta^{2}W>0$ for all
trial displacements~$\boldsymbol{\xi}$ satisfying~\eqref{eq:rel-7-46}
and the boundary conditions~\eqref{eq:rel-7-55}.  For axisymmetric
perturbations ($m=0$) this reduces directly to the requirement
$\Phi_{\mathrm{rel}}\ge0$, recovering the relativistic Rayleigh
criterion.

For non-axisymmetric perturbations the energy principle provides a
sharper sufficiency condition than the discriminant criterion alone.
Specifically, writing $\delta^{2}W = \delta^{2}W_{\mathrm{pot}}
+ \delta^{2}W_{\mathrm{kin}}$ and using the Schwarz inequality, one
arrives at:
\begin{equation}
\label{eq:rel-energy-suff}
\Phi_{\mathrm{rel}}(r) > \frac{m^{2}}{r^{2}}
  \max_{r}\!\left(\frac{4\,\Lf^{4}\,\Omega^{2}}
  {\Phi_{\mathrm{rel}}}\right)^{-1}
\end{equation}
as a sufficient condition for the stability of mode~$m$, which is
automatically satisfied when $\Phi_{\mathrm{rel}}$ is large and
positive (i.e.\ far from marginal stability).

%% ================================================================
\subsection*{(e) \emph{Causality constraint on perturbation modes}}
\label{sec:rel-7-67e}

A feature absent from the classical analysis is the requirement that
all perturbation phase velocities remain subluminal.  For a normal
mode of the form~\eqref{eq:rel-7-39}, the axial and azimuthal phase
speeds are, respectively,
\begin{equation}
\label{eq:rel-7-phase}
v_{\mathrm{ph},z} = -\frac{\operatorname{Re}(p)}{k},\qquad
v_{\mathrm{ph},\theta}
  = -\frac{\operatorname{Re}(p)\,r}{m}.
\end{equation}
Causality demands
\begin{equation}
\label{eq:rel-7-causality}
\boxed{%
|v_{\mathrm{ph},z}|\le c
\qquad\text{and}\qquad
\bigl|v_{\mathrm{ph},\theta}+V(r)\bigr|\le c,}
\end{equation}
where the second condition ensures that the perturbation, as seen in
the laboratory frame, does not exceed~$c$.

For stable modes ($p$ real), the Howard semicircle
theorem~\eqref{eq:rel-howard} already confines $p_r$ to the range
$[-m\Omega_{\max},\,-m\Omega_{\min}]$, so that
$|v_{\mathrm{ph},\theta}+V|\le r\,\Omega_{\max}+V<c$ whenever the
background flow itself is subluminal.  The axial causality bound
restricts the admissible wave numbers to
\begin{equation}
\label{eq:rel-7-kbound}
|k| \ge \frac{|m|\,\Omega_{\max}}{c}.
\end{equation}
Modes with $|k|$ below this threshold would require superluminal
axial phase speeds and are therefore excluded from the physical
spectrum.

For unstable modes ($p_i>0$), the group velocity
\begin{equation}
\label{eq:rel-7-group}
v_g = \frac{\partial\,\operatorname{Re}(p)}{\partial k}
\end{equation}
must also satisfy $|v_g|\le c$.  In the relativistic framework this is
guaranteed by the hyperbolic character of the underlying equations:
the linearised relativistic Euler system (being a first-order symmetric
hyperbolic system) has a well-posed initial-value problem with a finite
domain of dependence determined by the sound cones.

\medskip
\noindent\textbf{Summary.}\quad
The classical results of \S\,67 carry over to the relativistic regime
with two essential modifications:
\begin{enumerate}
\item The Rayleigh discriminant $\Phi(r)$ is replaced everywhere by
  its relativistic counterpart $\Phi_{\mathrm{rel}}(r)$, which
  involves the Lorentz factor~$\Lf$ of the background flow.
  Consequently the stability boundary shifts: flows that are
  classically stable may become unstable (and vice versa) at
  relativistic rotation speeds.
\item A new constraint---\emph{causality}---restricts the admissible
  perturbation spectrum, excluding modes whose phase or group velocity
  would exceed~$c$.  This constraint is automatically enforced by the
  hyperbolic structure of the relativistic equations.
\end{enumerate}
In the non-relativistic limit $V/c\to0$ all expressions reduce to
those of~\S\,67 of the original text.

%% ================================================================
\subsection*{(f) \emph{Astrophysical application: inviscid disk
  stability and the Rayleigh discriminant profile}}
\label{sec:rel-7-67f}

The inviscid stability analysis of this section has a direct
application to Keplerian accretion disks around compact objects
\citep{Seguin1975,FrankKingRaine2002}.  For geodesic circular orbits
in the Schwarzschild and Kerr metrics, the relativistic Rayleigh
discriminant $\Phi_{\mathrm{rel}}(r)$ can be evaluated exactly.  In
Schwarzschild geometry ($a=0$),
\begin{equation}
\label{eq:rel-7-Phi-Schw}
\Phi_{\mathrm{rel}}^{\mathrm{Schw}}(r)
= \Omega_{K}^{2}\!\left(1-\frac{6\,r_{g}}{r}\right),
\end{equation}
which vanishes at $r=6\,r_{g}$ (the ISCO) and is positive for all
$r>6\,r_{g}$.  For Kerr black holes with spin $a/M$,
\begin{equation}
\label{eq:rel-7-Phi-Kerr}
\Phi_{\mathrm{rel}}^{\mathrm{Kerr}}(r)
= \Omega_{K}^{2}\!\left(1-\frac{6\,r_{g}}{r}
  +\frac{8\,a\,r_{g}^{1/2}}{r^{3/2}}
  -\frac{3\,a^{2}}{r^{2}}\right),
\end{equation}
where $\Omega_{K}=(r^{3/2}+a\,r_{g}^{1/2})^{-1}(GM)^{1/2}$ is the
Kerr orbital frequency.

Figure~\ref{fig:Phi-rel-disk} plots
$\Phi_{\mathrm{rel}}(r)/\Omega_{K}^{2}(10\,r_{g})$ for spin parameters
$a/M=0,\,0.5,\,0.9,\,0.998$.  The discriminant is positive (stable)
outside the ISCO and passes through zero at the ISCO for each spin.
Higher spins allow stable orbits closer to the hole, in direct
correspondence with the narrower Couette gap that results from a
smaller inner boundary.  The Howard semicircle theorem of
\S\,\ref{sec:rel-7-67c} further constrains non-axisymmetric
perturbation growth rates to lie within the range of local angular
velocities, providing an upper bound on QPO growth in the stable
region of the disk.

\begin{figure}[htbp]
\centering
\includegraphics[width=0.75\textwidth]{fig_Phi_rel_schwarzschild_kerr.pdf}
\caption{Relativistic Rayleigh discriminant
  $\Phi_{\mathrm{rel}}(r)$ for Keplerian circular orbits in
  Schwarzschild and Kerr spacetimes.  The discriminant vanishes at the
  ISCO (filled circles) and is positive outside.  Higher black-hole
  spin shifts the ISCO inward, extending the stable Keplerian flow
  region.  The shaded area marks the Rayleigh-unstable regime
  ($\Phi_{\mathrm{rel}}<0$).}
\label{fig:Phi-rel-disk}
\end{figure}
